% questions.tex —— 2005 年·普通高等学校招生考试·重庆卷·文科 · 题目内容片段（无答案）
% 由 tools/generate_projects.py 依来源总账自 2005/2005_数学真题_重庆文科_试题.md 生成，勿手改（改源 MD 后重跑）

\gkpart{一、选择题（本大题共 10 小题，每小题 5 分，共 50 分）}

\begin{question}
\gkti{1}圆 $(x + 2)^2 + y^2 = 5$ 关于原点 $(0, 0)$ 对称的圆的方程为\mbox{（　　）}

\begin{gkopts}(2)
\task A．$(x - 2)^2 + y^2 = 5$
\task B．$x^2 + (y - 2)^2 = 5$
\task C．$(x + 2)^2 + (y + 2)^2 = 5$
\task D．$x^2 + (y + 2)^2 = 5$
\end{gkopts}

\end{question}
\begin{question}
\gkti{2}$\left(\cos\dfrac{\pi}{12} - \sin\dfrac{\pi}{12}\right)\left(\cos\dfrac{\pi}{12} + \sin\dfrac{\pi}{12}\right) =$\mbox{（　　）}

\begin{gkopts}(2)
\task A．$-\dfrac{\sqrt{3}}{2}$
\task B．$-\dfrac{1}{2}$
\task C．$\dfrac{1}{2}$
\task D．$\dfrac{\sqrt{3}}{2}$
\end{gkopts}

\end{question}
\begin{question}
\gkti{3}若函数 $f(x)$ 是定义在 $\mathbf{R}$ 上的偶函数，在 $(-\infty, 0]$ 上是减函数，且 $f(2) = 0$，则使得 $f(x) < 0$ 的 $x$ 的取值范围是\mbox{（　　）}

\begin{gkopts}(2)
\task A．$(-\infty, 2)$
\task B．$(2, +\infty)$
\task C．$(-\infty, -2) \cup (2, +\infty)$
\task D．$(-2, 2)$
\end{gkopts}

\end{question}
\begin{question}
\gkti{4}设向量 $\vec{a} = (-1, 2), \vec{b} = (2, -1)$，则 $(\vec{a} \cdot \vec{b})(\vec{a} + \vec{b})$ 等于\mbox{（　　）}

\begin{gkopts}(4)
\task A．$(1, 1)$
\task B．$(-4, -4)$
\task C．$-4$
\task D．$(-2, -2)$
\end{gkopts}

\end{question}
\begin{question}
\gkti{5}不等式组 $\begin{cases} |x - 2| < 2, \\ \log_2(x^2 - 1) > 1 \end{cases}$ 的解集为\mbox{（　　）}

\begin{gkopts}(4)
\task A．$(0, \sqrt{3})$
\task B．$(\sqrt{3}, 2)$
\task C．$(\sqrt{3}, 4)$
\task D．$(2, 4)$
\end{gkopts}

\end{question}
\begin{question}
\gkti{6}已知 $\alpha$、$\beta$ 均为锐角，若 $p : \sin\alpha < \sin(\alpha + \beta)$，$q : \alpha + \beta < \dfrac{\pi}{2}$，则 $p$ 是 $q$ 的\mbox{（　　）}

\begin{gkopts}(2)
\task A．充分而不必要条件
\task B．必要而不充分条件
\task C．充要条件
\task D．既不充分也不必要条件
\end{gkopts}

\end{question}
\begin{question}
\gkti{7}对于不重合的两个平面 $\alpha$ 与 $\beta$，给定下列条件：

① 存在平面 $\gamma$，使得 $\alpha$、$\beta$ 都垂直于 $\gamma$；

② 存在平面 $\gamma$，使得 $\alpha$、$\beta$ 都平行于 $\gamma$；

③ 存在直线 $l \subset \alpha$，直线 $m \subset \alpha$，使得 $l \parallel m$；

④ 存在异面直线 $l$、$m$，使得 $l \parallel \alpha$，$l \parallel \beta$，$m \parallel \alpha$，$m \parallel \beta$．

其中，可以判定 $\alpha$ 与 $\beta$ 平行的条件有\mbox{（　　）}

\begin{gkopts}(4)
\task A．1 个
\task B．2 个
\task C．3 个
\task D．4 个
\end{gkopts}

\end{question}
\begin{question}
\gkti{8}若 $(1 + 2x)^n$ 展开式中含 $x^3$ 项的系数等于含 $x$ 项的系数的 8 倍，则 $n$ 等于\mbox{（　　）}

\begin{gkopts}(4)
\task A．$5$
\task B．$7$
\task C．$9$
\task D．$11$
\end{gkopts}

\end{question}
\begin{question}
\gkti{9}若动点 $(x, y)$ 在曲线 $\dfrac{x^2}{4} + \dfrac{y^2}{b^2} = 1$（$b > 0$）上变化，则 $x^2 + 2y$ 的最大值为\mbox{（　　）}

\begin{gkopts}(1)
\task A．$\begin{cases} \dfrac{b^2}{4} + 4, & 0 < b < 4, \\ 2b, & b \geqslant 4 \end{cases}$
\task B．$\begin{cases} \dfrac{b^2}{4} + 4, & 0 < b < 2, \\ 2b, & b \geqslant 2 \end{cases}$
\task C．$\dfrac{b^2}{4} + 4$
\task D．$2b$
\end{gkopts}

\Needspace{0.4\textheight}
\end{question}
\begin{question}
\gkti{10}有一塔形几何体由若干个正方体构成，构成方式如图所示，上层正方体下底面的四个顶点是下层正方体上底面各连接中点，已知最底层正方体的棱长为 2，且该塔形的表面积（含最底层正方体的底面面积）超过 39，则该塔形中正方体的个数至少是\mbox{（　　）}
\nopagebreak[4]
\begin{center}\begin{tikzpicture}[x=1cm, y=1cm, line join=round, line cap=round, thick]
  \def\dx{0.45}
  \def\dy{0.38}

  % ================= Layer 1 (Bottom cube) =================
  \coordinate (A1) at (0, 0);
  \coordinate (B1) at (1.6, 0);
  \coordinate (C1) at ($(B1) + (\dx, \dy)$);
  \coordinate (D1) at ($(A1) + (\dx, \dy)$);

  \def\hOne{1.6}
  \coordinate (A1t) at ($(A1) + (0, \hOne)$);
  \coordinate (B1t) at ($(B1) + (0, \hOne)$);
  \coordinate (C1t) at ($(C1) + (0, \hOne)$);
  \coordinate (D1t) at ($(D1) + (0, \hOne)$);

  \filldraw[fill=white] (A1t) -- (B1t) -- (C1t) -- (D1t) -- cycle;
  \filldraw[fill=white] (A1) -- (B1) -- (B1t) -- (A1t) -- cycle;
  \filldraw[fill=white] (B1) -- (C1) -- (C1t) -- (B1t) -- cycle;

  % ================= Layer 2 (Middle cube) =================
  \coordinate (M1) at ($(A1t)!0.5!(B1t)$);
  \coordinate (M2) at ($(B1t)!0.5!(C1t)$);
  \coordinate (M3) at ($(C1t)!0.5!(D1t)$);
  \coordinate (M4) at ($(D1t)!0.5!(A1t)$);

  \def\hTwo{1.13}
  \coordinate (M1t) at ($(M1) + (0, \hTwo)$);
  \coordinate (M2t) at ($(M2) + (0, \hTwo)$);
  \coordinate (M3t) at ($(M3) + (0, \hTwo)$);
  \coordinate (M4t) at ($(M4) + (0, \hTwo)$);

  \filldraw[fill=white] (M1t) -- (M2t) -- (M3t) -- (M4t) -- cycle;
  \filldraw[fill=white] (M4) -- (M1) -- (M1t) -- (M4t) -- cycle;
  \filldraw[fill=white] (M1) -- (M2) -- (M2t) -- (M1t) -- cycle;

  % ================= Layer 3 (Top cube) =================
  \coordinate (N1) at ($(M4t)!0.5!(M1t)$);
  \coordinate (N2) at ($(M1t)!0.5!(M2t)$);
  \coordinate (N3) at ($(M2t)!0.5!(M3t)$);
  \coordinate (N4) at ($(M3t)!0.5!(M4t)$);

  \def\hThree{0.8}
  \coordinate (N1t) at ($(N1) + (0, \hThree)$);
  \coordinate (N2t) at ($(N2) + (0, \hThree)$);
  \coordinate (N3t) at ($(N3) + (0, \hThree)$);
  \coordinate (N4t) at ($(N4) + (0, \hThree)$);

  \filldraw[fill=white] (N1t) -- (N2t) -- (N3t) -- (N4t) -- cycle;
  \filldraw[fill=white] (N1) -- (N2) -- (N2t) -- (N1t) -- cycle;
  \filldraw[fill=white] (N2) -- (N3) -- (N3t) -- (N2t) -- cycle;
\end{tikzpicture}
\end{center}
\nopagebreak[4]
\begin{gkopts}(4)
\task A．$4$
\task B．$5$
\task C．$6$
\task D．$7$
\end{gkopts}

\end{question}
\gkpart{二、填空题（本大题共 6 小题，每小题 4 分，共 24 分）}

\begin{question}
\gkti{11}若集合 $A = \{x \in \mathbf{R} \mid x^2 - 4x + 3 < 0\}$，集合 $B = \{x \in \mathbf{R} \mid (x - 2)(x - 5) < 0\}$，则 $A \cap B =$ \underline{\hspace{2cm}}．

\end{question}
\begin{question}
\gkti{12}曲线 $y = x^3$ 在点 $(1, 1)$ 处的切线与 $x$ 轴、直线 $x = 2$ 所围成的三角形的面积为 \underline{\hspace{2cm}}．

\end{question}
\begin{question}
\gkti{13}已知 $\alpha$、$\beta$ 均为锐角，且 $\cos(\alpha + \beta) = \sin(\alpha - \beta)$，则 $\tan\alpha =$ \underline{\hspace{2cm}}．

\end{question}
\begin{question}
\gkti{14}若 $x^2 + y^2 = 4$，则 $x - y$ 的最大值是 \underline{\hspace{2cm}}．

\end{question}
\begin{question}
\gkti{15}若 10 把钥匙中只有 2 把能打开某锁，则从中任取 2 把能将该锁打开的概率为 \underline{\hspace{2cm}}．

\end{question}
\begin{question}
\gkti{16}已知 $A\left(-\dfrac{1}{2}, 0\right)$，$B$ 是圆 $F: \left(x - \dfrac{1}{2}\right)^2 + y^2 = 4$（$F$ 为圆心）上一动点，线段 $AB$ 的垂直平分线交 $BF$ 于 $P$，则动点 $P$ 的轨迹方程为 \underline{\hspace{2cm}}．

\end{question}
\gkpart{三、解答题（本大题共 6 小题，共 74 分）}

\begin{question}
\gkti{17}若函数 $f(x) = \dfrac{1 + \cos 2x}{2\sin\left(\dfrac{\pi}{2} - x\right)} + \sin x + a^2 \sin\left(x + \dfrac{\pi}{4}\right)$ 的最大值为 $\sqrt{2} + 3$，试确定常数 $a$ 的值．

\end{question}
\begin{question}
\gkti{18}加工某种零件需经过三道工序，设第一、二、三道工序的合格率分别为 $\dfrac{9}{10}$、$\dfrac{8}{9}$、$\dfrac{7}{8}$，且各道工序互不影响．

(1) 求该种零件的合格率；

(2) 从该种零件中任取 3 件，求恰好取到一件合格品的概率和至少取到一件合格品的概率．

\end{question}
\begin{question}
\gkti{19}设函数 $f(x) = 2x^3 - 3(a + 1)x^2 + 6ax + 8$，其中 $a \in \mathbf{R}$．

(1) 若 $f(x)$ 在 $x = 3$ 处取得极值，求常数 $a$ 的值；

(2) 若 $f(x)$ 在 $(-\infty, 0)$ 上为增函数，求 $a$ 的取值范围．

\Needspace{0.4\textheight}
\end{question}
\begin{question}
\gkti{20}如图，在四棱锥 $P - ABCD$ 中，底面 $ABCD$ 为矩形，$PD \perp$ 底面 $ABCD$，$E$ 是 $AB$ 上一点，$PE \perp EC$．已知 $PD = \sqrt{2}, CD = 2, AE = \dfrac{1}{2}$，求：
\nopagebreak[4]
\begin{center}\begin{tikzpicture}[>=stealth, scale=1.3, line join=round, line cap=round, thick]
  % Coordinates
  \coordinate (A) at (0, 0);
  \coordinate (D) at (0.8, 0.7);
  \coordinate (P) at (0.8, 2.8);
  \coordinate (E) at (1.4, 0);
  \coordinate (B) at (3.6, 0);
  \coordinate (C) at ($(B) + (D) - (A)$);

  % Dashed lines (hidden / interior)
  \draw[dashed] (P) -- (D);
  \draw[dashed] (D) -- (A);
  \draw[dashed] (D) -- (C);
  \draw[dashed] (D) -- (E);
  \draw[dashed] (E) -- (C);

  % Solid lines (outer / visible)
  \draw (P) -- (A);
  \draw (P) -- (E);
  \draw (P) -- (B);
  \draw (P) -- (C);
  \draw (A) -- (B) -- (C);

  % Labels
  \node[above] at (P) {$P$};
  \node[above right] at (D) {$D$};
  \node[below left] at (A) {$A$};
  \node[below] at (E) {$E$};
  \node[below right] at (B) {$B$};
  \node[above right] at (C) {$C$};
\end{tikzpicture}
\end{center}
\nopagebreak[4]
(1) 异面直线 $PD$ 与 $EC$ 的距离；

(2) 二面角 $E - PC - D$ 的大小．

\end{question}
\begin{question}
\gkti{21}已知中心在原点的双曲线 $C$ 的右焦点为 $(2, 0)$，右顶点为 $(\sqrt{3}, 0)$．

(1) 求双曲线 $C$ 的方程；

(2) 若直线 $l: y = kx + \sqrt{2}$ 与双曲线 $C$ 恒有两个不同的交点 $A$ 和 $B$，且 $\overrightarrow{OA} \cdot \overrightarrow{OB} > 2$（其中 $O$ 为原点）．求 $k$ 的取值范围．

\end{question}
\begin{question}
\gkti{22}数列 $\{a_n\}$ 满足 $a_1 = 1$ 且 $8a_{n+1} a_n - 16a_{n+1} + 2a_n + 5 = 0$（$n \geqslant 1$）．记 $b_n = \dfrac{1}{a_n - \dfrac{1}{2}}$（$n \geqslant 1$）．

(1) 求 $b_1$、$b_2$、$b_3$、$b_4$ 的值；

(2) 求数列 $\{b_n\}$ 的通项公式及数列 $\{a_n b_n\}$ 的前 $n$ 项和 $S_n$．

\end{question}